\section*{一般拓扑空间上的测度与狭收敛}
\phantomsection
\addcontentsline{toc}{section}{一般拓扑空间上的测度与狭收敛}

拓扑学与测度论之间联系的研究尤其被视为对测度的正则性性质的研究，特别是对“外”正则性与“内”正则性的研究；\(^{14}\) 在一个无穷可数的局部紧空间上，内正则性等价于外正则性。勒贝格给出的直线上集合的测度的构造凸显了这两种正则性，而波兰空间上测度的外正则性这一性质，似乎在 1935 年前后已是尽人皆知。但直到 1940 年，A. D. 亚历山德罗夫{[}3{]}才在一篇因战争而传播受阻的文章中凸显了内正则性的作用，并证明波兰空间上的测度具有这种性质；这一结果后来被普罗霍罗夫{[}254{]}重新发现，并且常被错误地归功于这位作者。直到很近的时候人们才看到这一性质可以推广到苏斯林空间；因此，这些空间的重要性大大提高了，尤其是因为人们了解到它们的理论可以不用可度量化假设，而且几乎全部函数空间都是苏斯林的（甚至多数情况下是卢津的）。\(^{15}\)

测度的收敛方式（弱的或狭的）的定义，最方便的做法是把测度空间与一个连续函数空间置于对偶之中。A. A. 马尔可夫推广 F. 里斯的一个古老结果，于 1938 年建立了 \(\mathcal{C}(X)\) 上的正泛函与紧空间 X 上的正则测度之间的一个双射对应。在前面已引用的著作{[}3{]}中，A. D. 亚历山德罗夫把这些结果推广到完全正则空间的情形：他在完全正则空间 X 上的有界连续函数空间 \(\mathcal{C}^{b}(X)\) 上的正线性形式的集合上引入一个层级，\(^{16}\) 他定义了有界测度的狭收敛，并特别证明了下面两个定理：

\begin{enumerate}
\def\labelenumi{\alph{enumi})}
\item
  若 X 是波兰空间，则 \(\mathcal{C}^{b}(X)\) 上与测度相对应的线性形式所成的集在序列的弱收敛下是闭的；
\item
  若有界测度的序列有一个狭极限，则“没有质量逃逸到无穷”（这是普罗霍罗夫狭收敛定理的逆定理的一个弱形式）。
\end{enumerate}

从这些概念与定理的融合中，普罗霍罗夫懂得如何为随机过程理论提取出一些重要结果，并以简单而引人注目的形式呈现它们。在他 1956 年那部已引用过的伟大著作{[}254{]}中，有相当一部分专门讨论波兰空间上的有界正测度；他推广莱维的一个构造，在质量为 1 的正测度所成的集合上定义了一个距离，使它成为一个波兰空间，然后建立了一个关于狭收敛的重要紧性判据。与普罗霍罗夫相独立，勒康{[}197{]}得到了关于测度狭收敛的若干紧性结果；他没有对他所考虑的空间作任何可度量化假设，而且他的结果在局部紧情形化归为迪厄多内先前的定理。

